\section{Постановка}

Цель проекта DRL v2 --- проверить, можно ли представить веса большой языковой
модели как дискретный язык маршрутов без заметной потери качества относительно
dense-базовой модели. В отличие от Route B v1, здесь приоритетом является не
минимальный размер, а минимальный разрыв по качеству на реальной текстовой модели.

Основной целевой объект --- Llama 3.3 70B Instruct. Основные целевые метрики:

\begin{itemize}[leftmargin=1.5em]
\item перплексия на полном WikiText-2;
\item поведение на кодовых задачах HumanEval;
\item latency и throughput reference/runtime-path;
\item потребление VRAM при surgery и eval.
\end{itemize}

\section{Метод}

Каждый вес $w$ кодируется не одним квантизованным числом, а маршрутом
переменной глубины

$$
\hat{w} = \rho \sum_{i=1}^{L} d_i s_i,
$$

где $d_i \in \{-1,0,+1\}$, $L \leq L_{\max}$, $L_{\max}=12$, $\rho$ ---
post-row scale, а $s_i$ --- элементы лестницы для данного семейства слоёв.

Фактически используется двухпроходная схема:

\begin{enumerate}[leftmargin=1.7em]
\item сначала вес кодируется как route переменной глубины;
\item затем из реально встретившихся route строится codebook, и каждый вес хранит
ID маршрута, а не сам маршрут.
\end{enumerate}

Рабочий рецепт калибровки на реальных весах оказался таким:

\begin{itemize}[leftmargin=1.5em]
\item построчная нормализация: $w_{norm} = w / \max |w_{row}|$;
\item geometric seed: $(1.0, 0.5, 0.25, \dots)$;
\item pinned top: верхняя ступень жёстко фиксируется в $1.0$;
\item 20 итераций coordinate-descent refinement.
\end{itemize}

Именно этот рецепт убрал ранний провал на реальных тензорах и перевёл relMSE в
диапазон порядка $10^{-5}$.